% !PRINTED 559
\weberpage{559}

%<<CONT
\[
\begin{aligned}
(a_{0,0} - \eta)\,\eta + &\ a_{1,0}\eta_1 + \cdots \hspace{2em} a_{e-1,0}\eta_{e-1} = 0 \\
a_{0,1}\eta \ + &\ (a_{1,1} - \eta)\eta_1 + \cdots \hspace{2em} a_{e-1,1}\eta_{e-1} = 0 \\
&\cdots\cdots\cdots\cdots\cdots\cdots\cdots\cdots\cdots\cdots \\
a_{0,e-1} \ + &\ a_{1,e-1}\eta_1 + \cdots (a_{e-1,e-1} - \eta)\eta_{e-1} = 0,
\end{aligned}
\]
% [sic] The printed last row omits the factor \eta after a_{0,e-1}; retained diplomatically although the parallel rows require it.
one sees, since the $\eta, \eta_1 \ldots \eta_{e-1}$ do not vanish, that the determinant of this system must be zero, hence

\weq{17}{\begin{vmatrix} a_{0,0}-\eta, & a_{1,0}, & \ldots & a_{e-1,0} \\ a_{0,1}, & a_{1,1}-\eta, & \ldots & a_{e-1,1} \\ \cdots & \cdots & \cdots & \cdots \\ a_{0,e-1}, & a_{1,e-1}, & \ldots & a_{e-1,e-1}-\eta \end{vmatrix} = 0,}

which, if one writes $x$ for $\eta$, must agree, apart from the sign, with $F_e(x)$.

\gesperrt{The function $F_e(x)$ therefore has not merely rational, but integer coefficients.}

% [sic] authority says that r is indeterminate here, although the displayed polynomial has indeterminate x and r is one of its roots.
One can now also form the equation of degree $f$ whose roots are $r, r_e, r_{2e} \ldots r_{(f-1)e}$; for, for an indeterminate $r$,

\weq{18}{\Phi_e(x) = (x - r)(x - r_e) \cdots (x - r_{(f-1)e})}

is a function in the field $R(\eta)$. Its coefficients can easily be computed from the \gesperrt{Newton} formulae, by which the coefficients of an equation are expressed by means of the power sums of the roots (\S.\,42). For by (12) and (3) we have
\[
r^{g^h} + r_e^{g^h} + \cdots + r_{(f-1)e}^{g^h} = \eta_h,
\]
and this is thus directly equal to one of the power sums, and every power sum is equal to one of the periods.

But the equation (18) can be treated in just the same way as the equation (2); for the group of $\Phi_e(x) = 0$ is precisely $C_e$, and if $e'$ is a divisor of $f$, say $f = e'f'$, then $C_{ee'}$ is a divisor of $C_e$. But to $C_{ee'}$ belongs the period
\[
\eta' = r + r_{ee'} + r_{2ee'} + \cdots + r_{(f'-1)ee'},
\]
which therefore depends on an equation of degree $e'$ in the field $R(\eta)$.

If one takes the numbers $e, e' \ldots$ to be primes, then the sequence of resolvents consists of a sequence of equations whose degrees are the prime factors of $n - 1$.

As for the group of the equation of degree $e$ whose roots are the $e$ quantities $\eta$, we find this very simply
%CONT>>
